\documentclass[12pt]{article}
\usepackage{amsmath}
\usepackage{graphicx}
\usepackage{geometry}
\geometry{a4paper, margin=1in}

\begin{document}

\begin{center}
    \textbf{University of Eswatini} \\
    \textbf{Department of Computer Science} \\
    \textbf{Examination (Main)} \\
    2023/2024 \\
    \textbf{FIRST SEMESTER} \\
    \vspace{1cm}
    \textbf{Title of Paper: INTRODUCTION TO COMPUTER SCIENCE} \\
    \textbf{Course Code: CSC111} \\
    \vspace{1cm}
    \textbf{Time Allowed: Three (3) Hours}
\end{center}

\noindent
\textbf{Instructions:} Answer \textbf{Question one} and any other \textbf{three} Questions.\\
Don’t write anything on the Examination Question paper.\\
You are not allowed to open this paper until you have been told to do so by the invigilator.

\vspace{1cm}

\noindent
\textbf{QUESTION ONE} (Compulsory)\\
a) List and explain the five distinct computing disciplines identified by ACM. \hfill (10 marks)\\
b) Define an algorithm and list five characteristics of a good algorithm. \hfill (5 marks)\\
c) Convert the binary number 101101 to decimal. \hfill (5 marks)\\
d) What is the role of the control unit in the Von Neumann architecture? \hfill (5 marks)\\

\vspace{0.5cm}

\noindent
\textbf{QUESTION TWO}\\
a) Describe the five generations of computers, highlighting the key technology and characteristics of each generation. \hfill (10 marks)\\
b) Differentiate between the following pairs:
    \begin{itemize}
        \item System software and application software
        \item Compiler and interpreter
        \item RAM and ROM
    \end{itemize}
    \hfill (9 marks)\\
c) Convert the hexadecimal number A3F to binary. \hfill (3 marks)\\
d) What is the stored program concept? \hfill (3 marks)\\

\vspace{0.5cm}

\noindent
\textbf{QUESTION THREE}\\
a) List and explain the six phases of the System Development Life Cycle (SDLC). \hfill (12 marks)\\
b) Explain why systems fail. List any five reasons. \hfill (5 marks)\\
c) What is the importance of documentation in system development? \hfill (5 marks)\\
d) Define the term "system" and list its key elements. \hfill (3 marks)\\

\vspace{0.5cm}

\noindent
\textbf{QUESTION FOUR}\\
a) Write an algorithm to calculate the area and circumference of a circle given the radius. \hfill (5 marks)\\
b) Draw a flowchart to represent the algorithm for finding the largest of three numbers. \hfill (10 marks)\\
c) Write pseudo-code for a program that asks the user for two numbers and prints their sum, difference, product, and quotient. \hfill (10 marks)\\

\vspace{0.5cm}

\noindent
\textbf{QUESTION FIVE}\\
a) Write a Python program to calculate the area of a triangle. The program should ask the user for the base and height and then display the area. \hfill (5 marks)\\
b) Write a Python program that takes a string from the user and does the following:
    \begin{itemize}
        \item Prints the string in uppercase.
        \item Prints the string in lowercase.
        \item Prints the length of the string.
        \item Prints the string with the first and last characters removed.
    \end{itemize}
    \hfill (10 marks)\\
c) Write a Python program that acts as a simple calculator. It should ask the user for two numbers and an operation (+, -, *, /) and then perform the operation and display the result. \hfill (10 marks)\\

\vspace{0.5cm}

\noindent
\textbf{QUESTION SIX}\\
a) Explain the following terms:
    \begin{itemize}
        \item Computer literacy
        \item Computer competency
        \item Information literacy
    \end{itemize}
    \hfill (6 marks)\\
b) Differentiate between a computer professional and a computer user, giving two examples of each. \hfill (5 marks)\\
c) List and explain five reasons why one should study computer science. \hfill (5 marks)\\
d) What are the side effects of using computers? List and explain any four. \hfill (5 marks)\\
e) Convert 2.5 GB to MB. \hfill (4 marks)\\

\end{document}